\documentclass[10pt,aspectratio=169]{beamer}

\usetheme{Singapore}
\setbeamertemplate{navigation symbols}{}

% Singapore's miniframes headline already shows the sections in colour and
% a dot per slide within each. All that is added here is a progress bar for
% position in the whole deck, and the frame number.
% the inactive sections are very faint by default; lift them enough to be
% read as "still to come" rather than as decoration
\setbeamercolor{section in head/foot}{fg=structure}
\setbeamercolor{section in head/foot shaded}{fg=black!45}
% Both arrows in one place, since a plain frame suppresses the footline
% and has to place them by hand.
\newcommand{\navcluster}{%
  \scriptsize
  \hyperlinkslideprev{\color{structure}\textless}\hspace{1.6ex}%
  \hyperlinkslidenext{\color{structure}\textgreater}%
}
\setbeamertemplate{footline}{%
  \hbox{\begin{beamercolorbox}[wd=\paperwidth,ht=1.6ex,dp=0.4ex,
      leftskip=1.4ex,rightskip=1.4ex]{}%
    % the frame number doubles as the one link back to the start, since
    % the navigation symbols are switched off above
    \navcluster\hfill\tiny\hyperlink{titleslide}{%
      \insertframenumber\,/\,\inserttotalframenumber}
  \end{beamercolorbox}}%
  \nointerlineskip
  % the pale bar is laid down at full width and the red one drawn over it,
  % rather than subtracting the two: the difference leaves a sub-point
  % sliver whenever the slide count does not divide the page width evenly
  \hbox to\paperwidth{%
    \rlap{{\color{bandpink}\rule{\paperwidth}{1.2pt}}}%
    {\color{structure}\rule{\dimexpr\paperwidth
      /\inserttotalframenumber*\insertframenumber\relax}{1.2pt}}%
    \hfill
  }%
}
\setbeamerfont{frametitle}{series=\bfseries}
\setbeamertemplate{itemize item}{--}
\setbeamertemplate{itemize subitem}{$\cdot$}

\usepackage{amsmath,amssymb,bm}
\usepackage{graphicx}
\usepackage[T1]{fontenc}

% Real Arial, which needs xelatex or lualatex; build with
%   latexmk -xelatex presentation.tex
% Maths stays in Latin Modern, as the reference deck also pairs Arial text
% with a serif maths font.
\usepackage{fontspec}
\setsansfont{Arial}
\renewcommand{\familydefault}{\sfdefault}
\usefonttheme{professionalfonts}

% The pink and red palette, sampled from
% 2026-08-29_10min_biophys_ML.pdf: a pale band behind the title, brick red
% for the title itself, brighter red for the current section.
\definecolor{bandpink}{RGB}{252,213,213}
\definecolor{deepred}{RGB}{168,48,36}
\definecolor{brightred}{RGB}{192,12,24}
\definecolor{mutedpink}{RGB}{226,150,150}

\setbeamercolor{structure}{fg=deepred}
\setbeamercolor{frametitle}{fg=deepred,bg=bandpink}
\setbeamercolor{title}{fg=deepred}
\setbeamercolor{subtitle}{fg=black}
\setbeamerfont{title}{size=\Huge,series=\mdseries}
\setbeamerfont{subtitle}{size=\large,series=\mdseries}
\setbeamercolor{normal text}{fg=black}
\setbeamercolor{palette primary}{fg=brightred,bg=bandpink}
\setbeamercolor{palette secondary}{fg=mutedpink,bg=bandpink!45}
\setbeamercolor{palette tertiary}{fg=mutedpink,bg=bandpink!45}
\setbeamercolor{section in head/foot}{fg=brightred,bg=bandpink}
\setbeamercolor{section in head/foot shaded}{fg=mutedpink,bg=bandpink!40}
\setbeamercolor{mini frame}{fg=brightred,bg=bandpink}
\setbeamercolor{item}{fg=deepred}
% a full-width band behind the title, as in the reference
\setbeamertemplate{frametitle}{%
  \nointerlineskip
  \begin{beamercolorbox}[wd=\paperwidth,ht=3.1ex,dp=1.3ex,
      leftskip=1.6ex,rightskip=1.6ex]{frametitle}%
    \centering\usebeamerfont{frametitle}\insertframetitle%
  \end{beamercolorbox}%
}

\newcommand{\Q}{\bm{Q}}
\newcommand{\pib}{\bm{\pi}}
\newcommand{\koff}{k_{\mathrm{off}}}
\newcommand{\code}[1]{\texttt{\small #1}}

\title{StatesAndRates}
\subtitle{Graphs of states connected by rates, and the linear algebra that
solves them}
\date{}

\begin{document}

% the title page gets a filled canvas, as in the reference deck
{
\setbeamercolor{background canvas}{bg=bandpink}
% not [plain]: that would drop the footline too, and with it the arrows.
% Suppressing only the headline keeps them where they are on every other
% slide, and keeps the section dots off the title. The head area is not
% covered by the background canvas, so an empty headline shows the paper
% through; it is filled here instead of left blank.
{\setbeamertemplate{headline}{\vbox to 0pt{\vss
    \hbox to\paperwidth{\color{bandpink}\rule{\paperwidth}{16pt}}}}
\begin{frame}
  \hypertarget{titleslide}{}%
  \vfill
  \begin{center}
    {\usebeamerfont{title}\usebeamercolor[fg]{title}\inserttitle\par}
    \vspace{1.0em}
    \begin{minipage}{0.9\textwidth}
      \centering\usebeamerfont{subtitle}\insertsubtitle\par
    \end{minipage}
    \vspace{2.2em}

    \begin{minipage}{0.72\textwidth}
      \centering\small
      A small package for building a kinetic mechanism once and asking it
      every question: equilibrium, relaxation, coarse graining, passage
      times, dwell times, single molecules.
    \end{minipage}

    \vspace{1.8em}
    {\small\url{https://github.com/lanmelo/StatesAndRates}}

  \end{center}
  \vfill
\end{frame}
}
}

\begin{frame}{A mechanism is a hypothesis}
  Not a summary of data --- a \textbf{claim}: which states exist, which
  interconvert, how fast.
  \[
    \underbrace{\text{states, arrows, rates}}_{\text{kinetic scheme}}
    \;\;\longleftrightarrow\;\;
    \underbrace{\text{wells, barriers, free energies}}_{\text{energy diagram}}
  \]
  \begin{center}
    \small one claim, drawn twice --- wherever detailed balance holds
  \end{center}
  \medskip
  So it predicts things an experiment can contradict:
  \begin{itemize}
    \item is a dissociation curve one exponential or two, and with what
      amplitudes?
    \item does a substitution move a well or a barrier --- association, or
      retention?
    \item what does an assay that cannot resolve the intermediate report as
      ``the'' off-rate?
    \item how often is an escape recaptured before it counts as one?
  \end{itemize}
  \medskip
  \begin{center}
    \textbf{Each is one operation on one matrix.}
  \end{center}
  \note{Whether the landscape exists at all is one more prediction ---
  section 2 asks when detailed balance holds. The four questions are the
  four middle sections of this deck. The package is those operations, so
  that the mechanism is the only thing written down by hand.}
\end{frame}

\section[Diagrams]{From diagrams to matrices}

\begin{frame}{The one object}
  A \code{Graph} is a set of states and the first-order rate constants
  between them. Everything else is a function of it.
  \begin{center}
    \includegraphics[width=0.92\textwidth]{figures/scheme.pdf}
  \end{center}
  \vspace{-0.6em}
  Rate constants are stored as \textbf{logarithms}, for three reasons:
  \begin{itemize}
    \item they span many orders of magnitude ($10^{-3}$ to $10^{2}\
      \mathrm{s^{-1}}$ here);
    \item a free-energy perturbation is then an \emph{additive} shift;
    \item an absent transition is $-\infty$, so structure and values live in
      one array.
  \end{itemize}
  A \code{Scheme} is the same topology with the rates \emph{named} rather
  than valued --- the mechanism before any numbers.
\end{frame}

\begin{frame}{From arrows to a matrix}
  A rate constant $k_{i\to j}$ is a \textbf{probability per unit time}: in a
  short interval $\mathrm{d}t$, a molecule in state $i$ moves to $j$ with
  probability $k_{i\to j}\,\mathrm{d}t$, regardless of how long it has
  already been in $i$.

  That single memoryless assumption fixes everything else:
  \begin{itemize}
    \item the time spent in $i$ before \emph{any} jump is exponentially
      distributed with rate $\sum_j k_{i\to j}$ --- the \emph{total} rate
      out;
    \item the destination is then chosen with probability
      $k_{i\to j}\big/\sum_{j'} k_{i\to j'}$, independently of the waiting
      time.
  \end{itemize}
  That is all ``continuous-time Markov chain'' means here. Nothing below
  needs more probability theory than this --- the rest is linear algebra on
  the matrix the diagram already is:
  \[
    \text{one arrow } i\to j
    \;\;\longleftrightarrow\;\;
    \text{one off-diagonal entry } k_{i\to j},
  \]
  with the diagonal fixed by bookkeeping, as the generator slide sets out.
\end{frame}

\begin{frame}{What first order buys, and what it costs}
  \textbf{Buys:} every rate is a constant, so the master equation on the
  next slide is \emph{linear}. That is what makes every question in this
  talk one operation on one matrix, with no time stepping and no iteration.

  \smallskip
  \textbf{Costs:} association is bimolecular --- a protein and a site have
  to find each other --- so it must enter as a \textbf{pseudo-first-order}
  rate with the concentration folded in:
  \[
    k_{\mathrm{on,max}}\ [\mathrm{s^{-1}}]
      \;=\; k^{(2)}_{\mathrm{on}}\ [\mathrm{M^{-1}s^{-1}}]
      \times [\mathrm{X}] .
  \]
  \begin{itemize}
    \item So every result belongs to \emph{that} concentration, not to a
      standard state. Free energies come out as $\Delta G$, not
      $\Delta G^\circ$.
    \item $k_{\mathrm{d}}$ is unaffected by concentration; $k_{\mathrm{a}}$
      scales with it. That asymmetry is why measured off-rates transfer
      between experiments and on-rates do not.
    \item Valid while the free pool is \emph{buffered},
      $[\mathrm{X}]\gg[\text{sites}]$. Depletion, competition and titration
      need a genuinely bimolecular treatment.
  \end{itemize}
\end{frame}

\begin{frame}{Notation, once}
  \begin{itemize}
    \item $\bm{p}(t)$ --- the \textbf{state vector}. Entry $p_i(t)$ is how
      much is in state $i$ at time $t$: a probability for one molecule, or a
      population for many. Either way the total is conserved, and it is what
      every time course in this talk plots.
    \item $\pib$ --- the \textbf{equilibrium distribution}: the
      $\bm{p}$ that has stopped changing. Entry $\pi_i$ is the
      equilibrium population of state $i$. Bold is the whole vector and a
      subscript picks out one entry, here and throughout: $\bm{p}$ and
      $p_i$, $\bm{K}$ and $K_{ij}$.
    \item $\bm{K}$ --- the \textbf{rate matrix} read straight off the
      diagram: $K_{ij}=k_{i\to j}$ off the diagonal, $K_{ii}=0$.
    \item $\bm{1}$ --- the \textbf{all-ones vector}. So $\bm{K}\bm{1}$ is the
      vector of \emph{row sums}, whose $i$th entry
      $\sum_j k_{i\to j}$ is the total rate out of state $i$ --- the
      quantity that already appeared as the exit rate on the last slide.
    \item $\mathrm{diag}(\bm{v})$ --- the matrix with $\bm{v}$ down its
      diagonal and zeros elsewhere.
  \end{itemize}
  \vspace{0.4em}
  Multiplying by $\bm{1}$ is the recurring trick: it sums along rows.
  ``$\bm{Q}\bm{1}=\bm{0}$'' will mean ``every row of $\bm{Q}$ sums to
  zero'', and the $-\bm{1}$ in the passage-time solve later is the constant
  on the right of one equation per state.
\end{frame}

\begin{frame}{The generator}
  Take the rate matrix and put minus each row's own sum on its diagonal:
  \[
    \Q = \bm{K} - \mathrm{diag}(\bm{K}\bm{1}),
    \qquad
    Q_{ij} = k_{i\to j}\ (i\neq j),
    \qquad
    Q_{ii} = -\sum_{j\neq i} k_{i\to j}.
  \]
  For the searching/testing/bound chain of slide~2, that is
  \[
    \Q =
    \begin{pmatrix}
      -k_{\mathrm{on,max}} & k_{\mathrm{on,max}} & 0\\[2pt]
      k_{\mathrm{off,M}} & -(k_{\mathrm{off,M}}+k_{\mathrm{on},\mu})
        & k_{\mathrm{on},\mu}\\[2pt]
      0 & k_{\mathrm{off},\mu} & -k_{\mathrm{off},\mu}
    \end{pmatrix} .
  \]
  \begin{itemize}
    \item Off-diagonal entries are the arrows. Zeros are the arrows that do
      not exist --- here searching and bound do not interconvert directly.
    \item Each diagonal entry is minus the rest of its row, so
      $\Q\bm{1}=\bm{0}$: \textbf{every row sums to zero}.
  \end{itemize}
\end{frame}

\begin{frame}{The master equation}
  Count what flows into a state and what flows out of it:
  \vspace{-0.4em}
  \[
    \frac{\mathrm{d}p_i}{\mathrm{d}t}
      = \underbrace{\sum_{j\neq i} k_{j\to i}\,p_j}_{
        \substack{\text{arriving from}\\ \text{every other state}}}
      - \underbrace{p_i\sum_{j\neq i} k_{i\to j}}_{
        \substack{\text{leaving }i\text{ by}\\ \text{every route out}}}
  \]
  That is the $i$th column of $\Q$ acting on $\bm{p}$, so the whole system
  is one matrix equation:
  \[
    \frac{\mathrm{d}\bm{p}}{\mathrm{d}t} = \Q^{\!\top}\bm{p} .
  \]
  \begin{itemize}
    \item The transpose is there because $Q_{ij}$ is the rate \emph{out of}
      $i$, while $\mathrm{d}p_i/\mathrm{d}t$ needs the rates \emph{into} it.
    \item Zero row sums make $\Q$ \textbf{singular} --- that is
      conservation, $\bm{1}^{\!\top}\mathrm{d}\bm{p}/\mathrm{d}t=0$.
    \item Everything from here is one linear-algebra question about $\Q$:
      a null vector, a matrix exponential, an eigenproblem, or a linear
      solve --- which one depends on what the experiment measures.
  \end{itemize}
\end{frame}

\section[Thermodynamics]{Thermodynamics}

\begin{frame}{Detailed balance: when do free energies exist?}
  Detailed balance is the statement that every transition is individually
  balanced at equilibrium,
  \[
    \pi_i\,k_{i\to j} = \pi_j\,k_{j\to i}
    \quad\Longleftrightarrow\quad
    \ln\frac{\pi_j}{\pi_i} = \ln\frac{k_{i\to j}}{k_{j\to i}} .
  \]
  Summing that around any cycle must give zero --- \textbf{Kolmogorov's
  criterion}. A graph that violates it is being driven and has no
  free-energy landscape.

  \textbf{How it is tested.} Not by enumerating cycles. Walk a spanning
  tree, accumulating $\ln\pi$ from the ratios; that fixes $\ln\pi$ up to a
  constant using $n-1$ of the transitions. Then the residual
  \[
    r_{ij} = \ln\pi_j - \ln\pi_i - \ln\frac{k_{i\to j}}{k_{j\to i}}
  \]
  must vanish on \emph{every remaining} transition. One tree walk,
  $O(n^2)$, no cycle basis.
\end{frame}

\begin{frame}{Equilibrium: a null vector, computed carefully}
  The equilibrium distribution is the stationary solution,
  \[
    \Q^{\!\top}\pib = \bm{0},
    \qquad \textstyle\sum_i \pi_i = 1 .
  \]
  $\Q$ is singular by construction, so forming this as a least-squares or
  nullspace problem invites catastrophic cancellation: the diagonal is a sum
  of the off-diagonals, and subtracting nearly equal large numbers destroys
  the small ones.

  \textbf{GTH reduction} instead eliminates one state at a time using
  \emph{only the non-negative off-diagonal rates}:
  \[
    Q_{ik} \leftarrow \frac{Q_{ik}}{\sum_{j<k} Q_{kj}},
    \qquad
    Q_{ij} \leftarrow Q_{ij} + Q_{ik}Q_{kj}
    \quad (i,j<k),
  \]
  then back-substitutes $w_k = \sum_{i<k} w_i Q_{ik}$ from $w_0=1$ and
  normalises. No subtraction ever occurs, so the result is accurate even
  when populations differ by many orders of magnitude.
\end{frame}

\begin{frame}{Free energies and energy diagrams}
  \begin{columns}[T]
    \begin{column}{0.5\textwidth}
      From the equilibrium distribution,
      \[ G_i = -RT\ln\pi_i , \]
      and from transition-state theory, writing
      $k_{i\to j} = A\,e^{-(G^\ddag_{ij}-G_i)/RT}$,
      \[ G^\ddag_{ij} = G_i - RT\ln\frac{k_{i\to j}}{A} . \]
      \begin{itemize}
        \item Both directions must give the \emph{same} barrier, which is
          detailed balance --- so the package refuses without it. The
          prefactor $A$ shifts every barrier equally, leaving all
          \emph{differences} unchanged.
      \end{itemize}
    \end{column}
    \begin{column}{0.5\textwidth}
      \includegraphics[width=0.78\textwidth]{figures/energy.pdf}
      \vspace{-0.7em}
      \begin{center}\scriptsize
        Same $k_{\mathrm{off},\mu}$, so the same well depth from the
        barrier; different recognition, different wells.

        \smallskip
        These are $\Delta G$, \emph{not} $\Delta G^\circ$: association
        enters as $k_{\mathrm{on}}[\mathrm{TF}]$, so this is drawn at that
        concentration. The standard-state shift is 12 kcal/mol.
      \end{center}
    \end{column}
  \end{columns}
\end{frame}

\begin{frame}{Perturbations are additive in log space}
  A \code{Mutation} is expressed in kcal/mol on the \emph{landscape}, not on
  individual rates, and applied as $\ln k \leftarrow \ln k + \Delta$:
  \[
    \underbrace{\Delta_{ij} \mathrel{+}= \frac{\Delta G_i}{RT}
      \ \ \forall j}_{\text{raise a well}}
    \qquad
    \underbrace{\Delta_{ij},\Delta_{ji} \mathrel{-}=
      \frac{\Delta G^\ddag_{ij}}{RT}}_{\text{raise a barrier}}
  \]
  \begin{itemize}
    \item Raising a well speeds up \emph{every} transition out of it ---
      one row of the matrix.
    \item Raising a barrier slows \emph{both} directions equally, so the
      equilibrium constant is untouched. Either way the graph stays
      thermodynamically consistent.
    \item $-\infty + \text{finite} = -\infty$: perturbing an absent
      transition does nothing, for free.
    \item Mutations add with \code{+}, so a double mutant is the sum of two
      single ones --- and that is testable against applying them in turn.
  \end{itemize}
\end{frame}

\section[Kinetics]{Kinetics}

\begin{frame}{Solving it: the matrix exponential}
  $\mathrm{d}\bm{p}/\mathrm{d}t = \Q^{\!\top}\bm{p}$ is the matrix version of
  $\dot{x}=ax$, and it has the same solution:
  \[
    \bm{p}(t) = e^{\Q^{\!\top} t}\,\bm{p}(0),
    \qquad
    e^{\bm{A}} \equiv \bm{I} + \bm{A} + \tfrac{1}{2}\bm{A}^2 + \dots
  \]
  So no time stepping is needed --- the answer at $t=10^{6}$ costs the same
  as at $t=1$.

  \textbf{But the series is not how you compute it.} Software squares a
  scaled exponential repeatedly, and each squaring doubles any error
  already present. That becomes visible once the largest rate times the
  longest time exceeds roughly $10^{9}$.

  \smallskip
  Master equations hit that constantly, because they are \textbf{stiff}
  whenever the timescales are well separated --- and well-separated
  timescales are usually the point. In the model here the fast mode is
  $10^{4}$ times the slow one, so following the slow relaxation means
  carrying the fast one for $10^{4}$ of its own lifetimes.

  \smallskip
  When the graph is reversible there is a way out that costs nothing.
\end{frame}

\begin{frame}{Solving it exactly: symmetrise first}
  $\Q$ is not symmetric, so its eigenvectors are not orthogonal and
  diagonalising it is ill-conditioned. Detailed balance repairs that.
  Conjugate by the square root of the equilibrium distribution,
  $\bm{\Pi}=\mathrm{diag}(\pib)$:
  \[
    \bm{S} = \bm{\Pi}^{1/2}\,\Q\,\bm{\Pi}^{-1/2},
    \qquad
    S_{ij} = \sqrt{\frac{\pi_i}{\pi_j}}\,k_{i\to j} .
  \]
  Now use $\pi_i k_{i\to j} = \pi_j k_{j\to i}$, which gives
  $\pi_i/\pi_j = k_{j\to i}/k_{i\to j}$:
  \[
    S_{ij} = \sqrt{\frac{k_{j\to i}}{k_{i\to j}}}\;k_{i\to j}
      = \sqrt{k_{i\to j}\,k_{j\to i}} \;=\; S_{ji}.
  \]
  \begin{itemize}
    \item \textbf{$\pib$ cancels.} The symmetrised matrix is just the
      geometric mean of each forward and reverse rate --- no equilibrium
      calculation needed to build it.
    \item It is symmetric, so \code{eigh} gives a real spectrum and genuinely
      orthogonal modes, at any stiffness.
  \end{itemize}
\end{frame}

\begin{frame}{Solving it exactly: putting it back}
  Conjugation does not change eigenvalues, so $\bm{S}$ and $\Q$ have the
  same spectrum. Diagonalise the easy one, $\bm{S}=\bm{V}\bm{\Lambda}
  \bm{V}^{\!\top}$, and undo the conjugation:
  \[
    \bm{p}(t) = \bm{\Pi}^{1/2}\,\bm{V}\,e^{\bm{\Lambda}t}\,
      \bm{V}^{\!\top}\,\bm{\Pi}^{-1/2}\,\bm{p}(0) .
  \]
  Read right to left, that is four steps: weight the start by
  $\bm{\Pi}^{-1/2}$, project onto the modes, decay each mode by its own
  $e^{\lambda t}$, then transform back.
  \begin{itemize}
    \item $e^{\bm{\Lambda}t}$ is diagonal --- one scalar exponential per
      mode --- so nothing is squared and nothing amplifies. It is exact
      however stiff the rates and however long the time.
    \item The cost is one eigendecomposition, reused for every requested
      time.
    \item Only the eigenvalues are needed if all you want is the
      \emph{rates} of the phases, which is the next slide.
  \end{itemize}
\end{frame}

\begin{frame}{The relaxation spectrum}
  \begin{columns}[T]
    \begin{column}{0.52\textwidth}\small
      The eigenvalues of $-\bm{S}$ are the relaxation rates:
      \[
        0 = \lambda_0 < \lambda_1 \le \dots \le \lambda_{n-1},
      \]
      $\lambda_0$ belonging to equilibrium and the rest being the rates of
      the exponential phases of \emph{any} relaxation. Their reciprocals are
      the relaxation times. For this chain,
      $\lambda_\pm = \tfrac{1}{2}\bigl(T \pm \sqrt{T^2-4D}\bigr)$,
      with $T=\sum_i k_i$ and
      $D = k_{\mathrm{off,M}}k_{\mathrm{off},\mu}
      + k_{\mathrm{off},\mu}k_{\mathrm{on,max}}
      + k_{\mathrm{on,max}}k_{\mathrm{on},\mu}$.
      \begin{itemize}
        \item A three-state chain has two non-zero modes, but here they are
          $10^{4}$ apart, so the slow one is all an experiment sees.
        \item That is why the population curve is indistinguishable from a
          single exponential with rate $k_{\mathrm{a}}+k_{\mathrm{d}}$ ---
          and it is the justification for coarse graining.
      \end{itemize}
    \end{column}
    \begin{column}{0.48\textwidth}
      \includegraphics[width=\textwidth]{figures/relaxation.pdf}
      \vspace{-0.4em}
      \begin{center}\small
        Populations from $\code{propagate}$; dashed is one exponential at
        $1/(k_{\mathrm{a}}+k_{\mathrm{d}})$.
      \end{center}
    \end{column}
  \end{columns}
\end{frame}

\begin{frame}{Absorbing states: the dissociation experiment}
  \begin{columns}[T]
    \begin{column}{0.5\textwidth}
      Setting one column of the rates to $-\infty$ in log space makes a
      state \textbf{absorbing}. Nothing rebinds, so the population only
      drains out of it.
      \begin{itemize}
        \item That is a wash step: protein removed from solution.
        \item The graph is no longer reversible, so \code{propagate} falls
          back to the matrix exponential --- correctly, since detailed
          balance no longer holds.
        \item Survival curves, and therefore measured off-rates, are
          exactly this calculation.
      \end{itemize}
    \end{column}
    \begin{column}{0.5\textwidth}
      \includegraphics[width=\textwidth]{figures/washout.pdf}
      \vspace{-0.4em}
      \begin{center}\small
        Association, then the searching state made absorbing at
        $t=2000\,$s.
      \end{center}
    \end{column}
  \end{columns}
\end{frame}

\section[Coarse graining]{What an experiment resolves}

\begin{frame}{Coarse graining: two different reductions}
  An experiment rarely resolves every state. Lumping states
  $\mathcal{A},\mathcal{B},\dots$ gives a graph over macrostates --- but
  there is more than one right answer, because the answer depends on what
  the instrument measures.

  \textbf{Local-equilibrium average} (\code{coarse\_grain}):
  \[
    K_{\mathcal{A}\mathcal{B}}
      = \sum_{i\in\mathcal{A}} w_i \sum_{j\in\mathcal{B}} k_{i\to j},
    \qquad
    w_i = \frac{\pi_i}{\sum_{i'\in\mathcal{A}}\pi_{i'}} .
  \]
  Reproduces the equilibrium populations and the equilibrium \emph{fluxes}
  exactly.

  \textbf{Passage-time inverse} (\code{passage\_rates}):
  \[
    K_{\mathcal{A}\mathcal{B}}
      = \Big(\sum_{i\in\mathcal{A}} w_i\, m_i(\mathcal{B})\Big)^{-1},
  \]
  the reciprocal of a mean first passage time. This is the rate a measured
  \emph{dwell time} reports.

  \smallskip
  They agree only when the partition is \textbf{lumpable}: every state of
  $\mathcal{A}$ has the same total rate into $\mathcal{B}$. Then the macro
  kinetics is genuinely Markovian.
\end{frame}

\begin{frame}{Mean first passage time: the first-step argument}
  How long, on average, until the system first reaches a target set
  $\mathcal{T}$? Let $m_i$ be that time starting from state $i$, with
  $m_i=0$ for $i\in\mathcal{T}$.

  Take one step and average. From $i$ you wait, then land somewhere:
  \[
    m_i \;=\;
    \underbrace{\frac{1}{\sum_j k_{i\to j}}}_{\substack{\text{mean wait}\\
      \text{in }i}}
    \;+\;
    \sum_j
    \underbrace{\frac{k_{i\to j}}{\sum_{j'} k_{i\to j'}}}_{\substack{
      \text{chance of}\\ \text{landing in }j}}
    \;\underbrace{m_j}_{\substack{\text{time still}\\ \text{to go}}} .
  \]
  Multiply through by the total exit rate and move the $m$ terms to one
  side. The waiting time and the branching probabilities recombine into the
  generator, and every trace of probability disappears:
  \[
    \boxed{\;\sum_{j\notin\mathcal{T}} Q_{ij}\, m_j = -1\;}
    \qquad (i\notin\mathcal{T}).
  \]
  One linear equation per non-target state, with the target states simply
  absent --- they contribute nothing because $m_j=0$ there.
\end{frame}

\begin{frame}{Mean first passage time: the linear solve}
  In matrix form, restricting to the non-target (``free'') states
  $\mathcal{F}$:
  \[
    \Q_{\mathcal{F}\mathcal{F}}\,\bm{m}_{\mathcal{F}} = -\bm{1}
    \qquad\Longrightarrow\qquad
    \bm{m}_{\mathcal{F}} = -\,\Q_{\mathcal{F}\mathcal{F}}^{-1}\,\bm{1} .
  \]
  With $\mathcal{F}=\{\text{testing},\text{bound}\}$, this is slide~7's
  matrix with the searching row and column struck out:
  \vspace{-0.9em}
  \[
    \Q_{\mathcal{F}\mathcal{F}} = \left(\begin{smallmatrix}
      -(k_{\mathrm{off,M}}+k_{\mathrm{on},\mu}) & k_{\mathrm{on},\mu}\\
      k_{\mathrm{off},\mu} & -k_{\mathrm{off},\mu}
    \end{smallmatrix}\right),
    \qquad
    \bm{m}_{\mathcal{F}} = \left(\begin{smallmatrix}
      k_{\mathrm{off},\mu}+k_{\mathrm{on},\mu}\\
      k_{\mathrm{off,M}}+k_{\mathrm{off},\mu}+k_{\mathrm{on},\mu}
    \end{smallmatrix}\right)
    \big/ k_{\mathrm{off,M}}k_{\mathrm{off},\mu} .
  \]
  \vspace{-0.9em}
  \begin{itemize}
    \item The full $\Q$ is singular, so this could not be solved as
      written on all states. \textbf{Deleting the target rows and columns
      is exactly what makes it invertible} --- the deleted rows are where
      the probability leaks out.
    \item $-\Q_{\mathcal{F}\mathcal{F}}^{-1}$ is the \textbf{fundamental
      matrix}: its $(i,j)$ entry is the expected \emph{total time spent in
      $j$} before absorption, starting from $i$. So $\bm{m}$ is just its
      row sums, and $-\Q_{\mathcal{F}\mathcal{F}}^{-1}\bm{1}$ says
      ``add up the time spent everywhere''.
    \item One solve, $O(n^3)$, gives the passage time from \emph{every}
      state at once --- not one target at a time.
    \item Everything downstream is this one inverse: the macroscopic rates
      of the next slide, and the observable dwell times later.
  \end{itemize}
\end{frame}

\begin{frame}{The same inverse, read two ways}
  The first-step argument gets to
  $\bm{m}=-\Q_{\mathcal{F}\mathcal{F}}^{-1}\bm{1}$ through probabilities.
  There is a route with no probability in it at all.

  The chance of still being free at time $t$, having started in $i$, is
  $\Phi_i(t)=\sum_{j\in\mathcal{F}}\big(e^{\Q_{\mathcal{F}\mathcal{F}}^{\!\top}t}
  \big)_{ji}$, and a mean lifetime is the integral of a survival curve:
  \[
    m_i = \int_0^\infty \Phi_i(t)\,\mathrm{d}t ,
    \qquad
    \int_0^\infty e^{\Q_{\mathcal{F}\mathcal{F}}t}\,\mathrm{d}t
      = -\,\Q_{\mathcal{F}\mathcal{F}}^{-1} ,
  \]
  the matrix analogue of $\int_0^\infty e^{-at}\mathrm{d}t = 1/a$. It
  converges precisely because absorption makes every eigenvalue negative.

  \smallskip
  So the fundamental matrix is the \textbf{time-integrated propagator}: its
  $(i,j)$ entry is the total time spent in $j$, the \emph{mean residence
  time}, and the passage time is the sum over $j$. That is the derivation in
  Bar-Haim \& Klafter, \emph{J.\ Chem.\ Phys.}\ \textbf{109}, 5187 (1998).

  \smallskip
  Both are available: \code{mean\_first\_passage\_times} for the totals in
  one solve, \code{mean\_residence\_times} for the matrix. Divide an entry
  by the mean time per visit and it becomes a \emph{visit count} ---
  $1/(1-p_{\mathrm{tot}})$ bound visits per escape, exactly.
\end{frame}

\begin{frame}{Recovering the macroscopic rates}
  \begin{columns}[T]
    \begin{column}{0.5\textwidth}
      For Marklund's model, lumping the invisible testing state and taking
      passage times between searching and bound reproduces the paper's
      Eq.~1 exactly:
      \[
        k_{\mathrm{a}} = k_{\mathrm{on,max}}\,p_{\mathrm{tot}},
        \qquad
        k_{\mathrm{d}} = k_{\mathrm{off},\mu}\,(1-p_{\mathrm{tot}}),
      \]
      with $p_{\mathrm{tot}} = k_{\mathrm{on},\mu}/
      (k_{\mathrm{on},\mu}+k_{\mathrm{off,M}})$ the recapture probability.
      \begin{itemize}
        \item No fitting, and no symbolic algebra either --- the package
          has none. This is the linear solve, which \emph{agrees} with
          Eq.~1 rather than deriving it.
        \item The three operators differ only in $p_{\mathrm{tot}}$, and
          they land on the predicted curve.
      \end{itemize}
    \end{column}
    \begin{column}{0.5\textwidth}
      \includegraphics[width=\textwidth]{figures/macro.pdf}
      \vspace{-0.4em}
      \begin{center}\small
        Line: Eq.~1 swept over $p_{\mathrm{tot}}$. Points: the three
        operators, from passage times.
      \end{center}
    \end{column}
  \end{columns}
\end{frame}

\begin{frame}{Splitting the passage time recovers Eq.~1 --- and its limits}
  \begin{center}
    \includegraphics[width=0.88\textwidth]{figures/residence.pdf}
  \end{center}
  \vspace{-0.8em}
  The residence matrix splits it:
  $m_{\mathrm{b}}
    = \underbrace{R_{\mathrm{bb}}}_{1/k_{\mathrm{d}}}
    + \underbrace{R_{\mathrm{bt}}}_{1/k_{\mathrm{off,M}}}$.
  There are $1/(1-p_{\mathrm{tot}})$ visits to bound, each lasting
  $1/k_{\mathrm{off},\mu}$, so their product is Marklund's
  $1/k_{\mathrm{d}}$ \textbf{exactly}. Testing is visited as often, but each
  visit lasts $1/(k_{\mathrm{on},\mu}+k_{\mathrm{off,M}})$ and the factors
  cancel, leaving one intermediate lifetime. So $1/\mathrm{MFPT}$ is Eq.~1
  only when that lifetime is short: the error is exactly $r/(1+r)$ with
  $r = k_{\mathrm{d}}/k_{\mathrm{off,M}}$.
\end{frame}

\section[Single molecules]{Single molecules}

\begin{frame}{Why simulate at all?}
  Everything so far computes what the population does \textbf{on average}.
  $\bm{p}(t)$ is a distribution over states, not a molecule.

  That is enough for a bulk time course, an equilibrium constant, or a
  relaxation rate. It is not enough when the observable is something the
  mean does not determine:
  \begin{itemize}
    \item the \emph{distribution} of dwell times, not just its mean --- is
      it one exponential or two?
    \item how many times a molecule revisits a state before it escapes;
    \item what a single-molecule trace looks like, which is what a
      microscope records;
    \item any statistic of a trajectory: waiting-time correlations, first
      passages beyond the mean, rare excursions.
  \end{itemize}
  A trajectory is one realisation of the \emph{same} generator. Nothing new
  is assumed --- and averaging many of them must reproduce
  \code{propagate}, which is a useful check on both.
\end{frame}

\begin{frame}{One step of the simulation}
  Every transition is first order, so molecules never interact: each is an
  independent copy of the same process and they advance together. Simulating 200 means 200 \emph{independent} copies, not
  200 molecules competing for one site.

  In state $i$, ask each outgoing arrow when it would fire:
  \vspace{-0.3em}
  \[
    \tau_{i\to j} = \frac{E_j}{k_{i\to j}},
    \qquad E_j \sim \mathrm{Exp}(1)\ \text{drawn independently},
  \]
  an exponential waiting time of rate $k_{i\to j}$, which is what that
  rate constant means. Then take the \textbf{first} arrow to fire:
  \[
    j^\star = \arg\min_j \tau_{i\to j},
    \qquad
    t \leftarrow t + \min_j \tau_{i\to j},
    \qquad
    i \leftarrow j^\star .
  \]
  Repeat --- hence \emph{first-reaction} method.

  \smallskip
  \textbf{The method is not the restriction.} Gillespie's algorithm is exact
  for any well-mixed mass-action network, bimolecular included: for
  $\mathrm{A}+\mathrm{B}\to\mathrm{C}$ the propensity is $c\,N_A N_B$,
  recomputed from the current counts each step. It is \emph{this package}
  that is first order, which is what lets it track independent molecules on
  a fixed graph rather than a vector of counts.
\end{frame}

\begin{frame}{Why that is the right answer}
  Two facts make the shortcut legitimate:
  \begin{itemize}
    \item the minimum of independent exponentials is exponential with the
      \emph{summed} rate, $\min_j \tau_{i\to j}\sim
      \mathrm{Exp}\!\big(\sum_j k_{i\to j}\big)$ --- which is the exit time
      the memoryless assumption demanded;
    \item the probability that arrow $j$ is the one that fires first is
      $k_{i\to j}\big/\sum_{j'} k_{i\to j'}$ --- which is the branching
      probability.
  \end{itemize}
  So drawing a race between the arrows reproduces ``sample the exit time,
  then choose a destination'' without ever writing those two steps
  separately.

  \smallskip
  \textbf{Done in log space.} $\ln\tau_{i\to j} = \ln E_j - \ln k_{i\to j}$,
  and the rates are already stored as logarithms. Rates spanning many orders
  of magnitude therefore cost no accuracy, and an absent transition is
  $\ln k = -\infty$, giving $\tau=+\infty$ rather than a division by zero.

  \smallskip
  An \textbf{absorbing} state has no finite $\tau$ at all, so the molecule
  simply stops --- no special case needed.
\end{frame}

\begin{frame}{What a single molecule looks like}
  \begin{center}
    \includegraphics[width=0.62\textwidth]{figures/gillespie.pdf}
  \end{center}
  \vspace{-0.4em}
  \begin{itemize}
    \item Same mechanism, two operators. Osym sits bound; O2 keeps
      returning to the search.
    \item The trajectory carries every jump, so occupancy at arbitrary
      times is a \code{searchsorted} plus a cumulative sum --- no
      re-simulation.
    \item The check on all the matrix algebra: the ensemble average of
      many trajectories reproduces \code{propagate}.
  \end{itemize}
\end{frame}

\begin{frame}{The ensemble average is the master equation}
  \begin{columns}[T]
    \begin{column}{0.54\textwidth}
      \includegraphics[width=\textwidth]{figures/aggregate.pdf}
    \end{column}
    \begin{column}{0.46\textwidth}
      \small Solid: 200 molecules through \code{occupancy}; dashed:
      \code{propagate} on the same graph.
      \begin{itemize}
        \item This is the one place the two halves of the package have to
          agree, and they share no code: NumPy sampling of waiting times
          against a JAX eigendecomposition.
        \item The residual is sampling noise, not bias: rms $0.026$
          against the binomial $1/\sqrt{N}=0.071$, and it shrinks with
          more molecules, not with a finer grid.
        \item One molecule at a time and a whole population measure the
          same object; which to run is a question about noise.
      \end{itemize}
    \end{column}
  \end{columns}
\end{frame}

\begin{frame}{Dwell times: two rates, and they are not the same}
  A dwell in one state can report either of two rates.

  \textbf{Per visit} --- arrival until the next jump out, whatever the
  destination:
  \[
    \text{exponential with}\quad k_{\text{micro}} = \sum_j k_{i\to j} .
  \]
  \textbf{Until it actually leaves} --- arrival until first reaching a
  target, so excursions that come straight back \emph{do not} end the dwell:
  \[
    k_{\text{obs}} = 1/m_i(\mathcal{T}) .
  \]
  \begin{itemize}
    \item The second is always slower, by one over the probability that an
      escape is \emph{not} recaptured.
    \item Which one an experiment measures depends on whether it can see the
      excursion. If the intermediate is invisible, it measures the passage
      time --- so the microscopic rate is not what a dwell histogram gives.
    \item This is the microscopic/macroscopic distinction, in one state.
  \end{itemize}
\end{frame}

\begin{frame}{Dwells, sampled and predicted}
  \begin{center}
    \includegraphics[width=0.5\textwidth]{figures/dwells.pdf}
  \end{center}
  \vspace{-0.6em}
  \begin{itemize}
    \item Histogram: Osym dwells extracted from a Gillespie run. Lines:
      $k\,e^{-kt}$ at the two rates above, computed from the graph, not
      fitted to the sample.
    \item The sample follows the \emph{observable} rate, not the
      microscopic one --- $1/k_{\mathrm{off},\mu}=167\,$s against a measured
      mean of $\sim\!1100\,$s for Osym.
    \item Dwells still running at the end are dropped, pulling the
      sample mean slightly low.
  \end{itemize}
\end{frame}

\section[Driven, differentiable]{Beyond closed form}

\begin{frame}{Time-dependent rates, and derivatives}
  \textbf{When the rates move.} A ligand jump, a flow ramp, a wash with
  finite mixing: then $\Q=\Q(t)$ and there is no closed form. The package
  integrates
  \[
    \frac{\mathrm{d}\bm{p}}{\mathrm{d}t} = \Q(t)^{\!\top}\bm{p}
  \]
  with an implicit solver (Kvaerno5), because master equations are stiff
  whenever their timescales are well separated --- which is whenever the
  problem is interesting.

  \textbf{Why it is all differentiable.} A graph is registered as a JAX
  pytree whose \emph{only} leaf is the log-rate matrix; the state names and
  temperature are static metadata. So
  \[
    \frac{\partial}{\partial \ln k_{i\to j}}
      \big(\text{any observable}\big)
  \]
  comes from \code{jax.grad}, and a library of variants is one
  \code{jax.vmap}. Fitting a mechanism to data is then the same machinery as
  evaluating it.
\end{frame}

\begin{frame}{Everything above, as matrix algebra}
  \small
  One operator, two directions. \textbf{Forward} questions act on a
  distribution and need the rates \emph{into} a state, so they carry
  $\Q^{\!\top}$; \textbf{backward} questions act on a function of where you
  \emph{started} and need the rates \emph{out}, so they carry $\Q$, with the
  target's rows and columns deleted.

  \vspace{0.1em}
  {\renewcommand{\arraystretch}{1.15}
  \begin{tabular}{@{}lllll@{}}
    & \textbf{Unknown} & \textbf{Equation} & \textbf{Name}
      & \textbf{What is measured}\\[2pt]
    \hline\\[-6pt]
    \textbf{Forward}
      & $\bm{\pi}$ & $\Q^{\!\top}\bm{\pi}=\bm{0}$,\;
        $\bm{1}^{\!\top}\bm{\pi}=1$ & stationary distribution
        & occupancy, $K_{\mathrm{d}}$, $\Delta G$\\
      & $\bm{p}(t)$ & $\bm{p}(t)=e^{\Q^{\!\top}t}\bm{p}(0)$ & propagation
        & a kinetic trace\\
      & $\bm{J}$ & $J_{ij}=\pi_i k_{i\to j}-\pi_j k_{j\to i}$
        & stationary flux & turnover, dissipation\\[3pt]
    \textbf{Backward}
      & $\bm{N}$ & $\Q_{\mathcal{F}\mathcal{F}}\bm{N}=-\bm{I}$
        & fundamental matrix & where the time goes\\
      & $\bm{m}$ & $\Q_{\mathcal{F}\mathcal{F}}\bm{m}=-\bm{1}$,\;
        $\bm{m}=\bm{N}\bm{1}$ & MFPT
        & dwell time, $1/\koff$\\
      & $\bm{H}$ & $\Q_{\mathcal{F}\mathcal{F}}\bm{H}
        =-\Q_{\mathcal{F}\mathcal{A}}$ & splitting probabilities
        & branching ratio\\[3pt]
    \textbf{Either} & $\lambda$ & $\Q\bm{v}=\lambda\bm{v}$
      & relaxation spectrum & exponential phases\\
  \end{tabular}}

  \vspace{0.4em}
  The two halves are one object read twice:
  $-\Q_{\mathcal{F}\mathcal{F}}^{-1}
    = \int_0^\infty e^{\Q_{\mathcal{F}\mathcal{F}}t}\,\mathrm{d}t$, and
  $e^{\Q^{\!\top}t}\to\bm{\pi}\bm{1}^{\!\top}$ as $t\to\infty$. So a
  backward solve is a forward propagation with the clock integrated out ---
  no $t$ in it, and every starting state at once --- and equilibrium is
  propagation at $t\to\infty$, only cheaper.
\end{frame}

\begin{frame}{The map}
  \scriptsize
  \begin{tabular}{@{}lll@{}}
    \textbf{Question} & \textbf{Linear algebra} & \textbf{Entry point}\\[2pt]
    \hline\\[-6pt]
    Is there a landscape? & tree walk $+$ residual &
      \code{is\_detailed\_balanced}\\
    Equilibrium populations & null vector of $\Q^{\!\top}$, by GTH &
      \code{equilibrium}\\
    Free energies, barriers & $-RT\ln\pi$; $G_i-RT\ln(k/A)$ &
      \code{free\_energies}, \code{barrier\_energies}\\
    Perturb the landscape & additive shift of $\ln k$ &
      \code{Mutation}, \code{Graph.mutate}\\
    Populations over time & $e^{\Q^{\!\top}t}$, or $\bm{\Pi}^{1/2}$-conjugated
      \code{eigh} & \code{propagate}\\
    Observable phases & eigenvalues of $-\bm{S}$ &
      \code{relaxation\_rates}, \code{\dots times}\\
    Driven rates & implicit ODE solve & \code{integrate}\\
    What a lumped experiment sees & flux average, or $1/$MFPT &
      \code{coarse\_grain}, \code{passage\_rates}\\
    Time to reach a state & $-\Q_{\mathcal{F}\mathcal{F}}^{-1}\bm{1}$ &
      \code{mean\_first\_passage\_times}\\
    Where that time is spent & $-\Q_{\mathcal{F}\mathcal{F}}^{-1}$ itself &
      \code{mean\_residence\_times}\\
    Is lumping legitimate? & equal row sums into each block &
      \code{is\_lumpable}\\
    Individual molecules & first-reaction sampling &
      \code{simulate}\\
    Dwell distributions & row sum vs.\ MFPT &
      \code{Trajectory.dwells}, \code{dwell\_rates}\\
    Draw it & --- & \code{state\_diagram}, \code{energy\_diagram}\\
  \end{tabular}

  \vspace{0.8em}
  \normalsize
  Build the mechanism once as a \code{Scheme}; fill it with numbers to get a
  \code{Graph}; ask any of the above.
\end{frame}

\end{document}
